\section{Introduction}
\label{sec:intro}

Machine learning (ML) has become a standard tool of empirical finance \cite{gu2020empirical}, and a common way to
use it in trading is to attach a learned model to a rule-based strategy that is already understood and validated
\cite{lopezdeprado2018afml,joubert2022metalabeling}. The evidence offered for such hybrid systems is usually a single
comparison: the hybrid against the rule-based baseline, summarized by one difference in Sharpe ratio. That comparison
answers whether the hybrid performed better in one backtest. It does not answer the question a designer actually faces,
namely \emph{at which decision} a learned model should replace the rule.

A hybrid pipeline has several natural insertion points: whether a bar is admissible for trading (a regime filter),
whether a fired signal should be taken (meta-labeling), how much to risk on it (position sizing), when to enter, and
when to exit. An aggregate before-and-after comparison cannot attribute an improvement or a loss to any one of them.
It is also confounded in ways that are easy to overlook: each learned component typically receives its own features,
tuning budget, and evaluation window, so a difference between two hybrids can reflect feature engineering or
evaluation leniency rather than the insertion point. Finally, the literature on technical trading rules shows that
apparent performance is fragile to the number of configurations searched
\cite{sullivan1999data,white2000reality,harvey2016cross}; reporting the one hybrid that looked best compounds that
problem.

This paper introduces \emph{matched-budget component attribution} (MBCA), a controlled-variable protocol for this
question. MBCA holds the data, the feature set, the model class and its hyperparameters, the freeze discipline, and
the evaluation protocol fixed, and varies only which single decision of a validated rule-based system a learned model
replaces. Because everything else is matched, a change in performance can be attributed to the insertion point
itself, and a null result is a precise statement about the stated budget at that insertion point.

We apply MBCA twice. \emph{Study~1} applies it to one validated strategy, a Jurik moving-average (JMA) crossover, in
three markets and finds \SoneSigGain{} significant improvements and \SoneSigLoss{} significant losses among
\SoneConfigs{} configurations; an independent re-implementation reproduces this headline but not its market-level
claims. \emph{Study~2} is a pre-registered test \cite{nosek2018preregistration} of whether Study~1 generalizes. Its
design was committed to version control before any of its data existed, and it evaluates five historically documented
strategies on \STwoInstruments{} daily series across \STwoMarkets{} asset-class markets, at two independent freeze
dates, under conservative execution, with mark-to-market portfolio accounting and a paired block-bootstrap test with
false-discovery-rate (FDR) control.

The paper makes four contributions:
\begin{enumerate}
  \item \emph{Method.} MBCA and a taxonomy of insertion points (independent, narrowing, replacing) that predicts how
        much a component can change a rule-based system.
  \item \emph{Evidence.} \STwoMktTests{} pre-registered market-level tests of five insertion points and one
        combination across five strategy families, with replication at a second freeze and a battery of
        generalization, robustness, and bias checks.
  \item \emph{Auditability.} One configuration file drives data collection, every experiment, and every table,
        figure, and in-text number; a test suite enforces the methodological invariants and fails if the manuscript
        contains a hand-typed result.
  \item \emph{Software.} An open-source implementation with which practitioners can apply MBCA to their own strategy
        and data.
\end{enumerate}

The evidence can be summarized in five findings. First, once multiplicity is controlled, ML rarely changes a
historically successful rule significantly: \STwoMktSigGain{} significant gains and \STwoMktSigLoss{} significant
losses in \STwoMktTests{} tests. Second, narrowing components (meta-labeling and position sizing) are the only ones
that are directionally positive in most cells and never significantly harmful, and their effect transfers across
markets; it is, however, small and replicates only partially across freezes. Third, replacing components are strategy
dependent: ML entry is the most harmful component, whereas ML exit, which is negative on average for
trend-followers, is the only component that significantly improves a strategy in several markets, in each case for
RSI(2) mean reversion. Fourth, execution assumptions dominate: filling breached stops at the stop level rather than at
the breaching close inflates the controls' Sharpe ratios by \STwoExecBaseStopMin{} to \STwoExecBaseStopMax{} and made
ML exit look worse in Study~1 than it is. Fifth, the controls themselves lose most of their in-sample edge out of
sample, which bounds what any component can add.

The remainder of the paper is organized as follows. Section~\ref{sec:related} reviews related work.
Section~\ref{sec:method} defines MBCA and its taxonomy. Section~\ref{sec:study1} reports Study~1.
Sections~\ref{sec:study2design} and~\ref{sec:study2results} present the design and results of Study~2.
Section~\ref{sec:discussion} discusses the findings and their practical implications, Section~\ref{sec:limitations}
states the limitations, and Section~\ref{sec:conclusion} concludes.
